\chapter{The Gradient}

\epigraph{\textit{For mostly they goes up and down \ldots}}{P.~R.~Chalmers}

%% ================================================================
\section{Line Integrals and the Gradient}
%% ================================================================

We have now investigated the relationship between the following
two statements:
\begin{enumerate}
\item $\oint_C \bfF \cdot \hat{\mathbf{t}}\ds = 0$ for any closed curve $C$.
\item $\nabla \times \bfF = 0$.
\end{enumerate}

We saw in the last chapter that the first of these statements implies
the second and is equivalent to the assertion that the line integral
of $\bfF \cdot \hat{\mathbf{t}}$ is independent of the path.  We also saw that the
second statement implies the first if $\bfF$ is smooth in a simply
connected region.  You might think that two ways of saying
something would be enough, but there is a \emph{third} way, as we
shall now see.

Let us suppose that a given vector function $\bfF(x,\,y,\,z)$ has associated
with it a scalar function $\psi(x,\,y,\,z)$ and that the two functions
are related as follows:
\begin{equation}
F_x = \pd{\psi}{x}\,,\qquad F_y = \pd{\psi}{y}\,,\qquad\text{and}\qquad
F_z = \pd{\psi}{z}\,.
\label{eq:IV-1}
\end{equation}

If the preceding relations hold, then the line integral of
$\bfF \cdot \hat{\mathbf{t}}$ is independent of path.  To show this, we use the
three relations given in Equation~(\ref{eq:IV-1}) and the formula
for the unit tangent vector to get
\[
\bfF \cdot \hat{\mathbf{t}}
= \pd{\psi}{x}\frac{dx}{ds}
+ \pd{\psi}{y}\frac{dy}{ds}
+ \pd{\psi}{z}\frac{dz}{ds}
= \frac{d\psi}{ds}
\]
where the second equality follows from a familiar chain rule of
multivariate calculus.  Suppose now that the path $C$ joins the two
points $(x_0,\,y_0,\,z_0)$ and $(x_1,\,y_1,\,z_1)$.  Then
\[
\int_C \bfF \cdot \hat{\mathbf{t}}\ds
= \int_C \frac{d\psi}{ds}\ds
= \int_C d\psi
= \psi(x_1,\,y_1,\,z_1) - \psi(x_0,\,y_0,\,z_0).
\]

You can see that this result depends only on the points at which the
path $C$ begins and ends.  We'd get the same result for \emph{any}
path joining these two points.  This proves our assertion: with
$\bfF$ and $\psi$ related as in Equations~(\ref{eq:IV-1}), the line
integral of $\bfF \cdot \hat{\mathbf{t}}$ is independent of path.  We shall
now show that the converse of this statement is also true; that is, if
the line integral of $\bfF \cdot \hat{\mathbf{t}}$ is independent of path, there
is a scalar function $\psi(x,\,y,\,z)$ related to $\bfF$ as specified in
Equations~(\ref{eq:IV-1}).

We begin with the observation that, because the line integral
$\int_C \bfF \cdot \hat{\mathbf{t}}\ds$ is independent of path, if we integrate
from some \emph{fixed} point $P_0(x_0,\,y_0,\,z_0)$ to a second point
$P(x,\,y,\,z)$, the result is a scalar function of the coordinates
$(x,\,y,\,z)$:
\begin{equation}
\psi(x,\,y,\,z)
= \int_{(x_0,y_0,z_0)}^{(x,y,z)} \bfF \cdot \hat{\mathbf{t}}\ds.
\label{eq:IV-2}
\end{equation}

It is important to understand that this would not be true if the
integral depended on path, for then its value would depend not only
on the coordinates $(x,\,y,\,z)$ of the point $P$ but also on the path
joining $P_0$ and $P$, and the integral would not then be a function
within the standard definition of the term.

Since the integral we're examining is path independent, we are
free to select any curve as the path of integration.  We choose the
one shown in Figure~IV--1.  It consists of two parts.  The first, $C_0$,

\begin{figure}[ht]
\centering
\includegraphics[width=0.6\textwidth]{figures/ch04/fig_IV_1.png}
\caption{}
\label{fig:IV-1}
\end{figure}

connects $P_0$ to an intermediate point $P_1$ whose coordinates are
$(a,\,y,\,z)$, where $a$ is some constant.  Beyond fixing its two end
points and requiring it to be reasonably smooth, we do not need to
specify anything more about $C_0$.  The second part of the curve,
$C_1$, is the straight-line segment from $P_1$ to $P$.  Thus,
Equation~(\ref{eq:IV-2}) becomes
\[
\psi(x,\,y,\,z)
= \int_{P_0}^{P_1} \bfF \cdot \hat{\mathbf{t}}\ds
+ \int_{P_1}^{P} F_x(x',\,y,\,z)\,dx'.
\]

The first term on the right-hand side of this equation is independent
of the variable $x$.  The second term is, effectively, nothing more
than an ordinary one-dimensional integral, since $y$ and $z$ are
constant on $C_1$ and just come along for the ride.  That is,
\[
\int_{P_1}^{P} F_x(x',\,y,\,z)\,dx'
= \int_a^x F_x(x',\,y=\text{const.},\,z=\text{const.})\,dx',
\]
and so
\begin{align*}
\pd{\psi}{x}
&= \frac{d}{dx}\int_a^x F_x(x',\,y=\text{const.},\,z=\text{const.})\,dx'\\
&= F_x(x,\,y,\,z),
\end{align*}
where we use the fact that the derivative of an integral with respect
to its upper limit is merely the integrand evaluated at that

\begin{figure}[ht]
\centering
\includegraphics[width=0.5\textwidth]{figures/ch04/fig_IV_2a.png}
\caption{}
\label{fig:IV-2a}
\end{figure}

\begin{figure}[ht]
\centering
\includegraphics[width=0.5\textwidth]{figures/ch04/fig_IV_2b.png}
\caption{}
\label{fig:IV-2b}
\end{figure}

\noindent limit.  This establishes one of the three relations we
sought.  The other two, $F_y = \partial\psi/\partial y$ and
$F_z = \partial\psi/\partial z$, can be obtained by the same sort of
reasoning, and you should carry out the derivations yourself.
Figure~IV--2(a) and (b) will be helpful.

You have probably recognized by now that we have here another
use for the del notation.  That is,
\[
F_x = \pd{\psi}{x}\,,\qquad F_y = \pd{\psi}{y}\,,\qquad\text{and}\qquad
F_z = \pd{\psi}{z}
\]
can be combined to give
\[
\bfF = \bfi\,\pd{\psi}{x} + \bfj\,\pd{\psi}{y} + \bfk\,\pd{\psi}{z}
= \left(\bfi\,\frac{\partial}{\partial x}
  + \bfj\,\frac{\partial}{\partial y}
  + \bfk\,\frac{\partial}{\partial z}\right)\psi
= \nabla\psi,
\]
which is read ``del psi.''  This operator is called the \emph{gradient}
and is sometimes written $\Grad\psi$.  However, we shall always
write $\nabla\psi$ in keeping with modern usage.  The gradient of
$\psi$ is a \emph{vector function} of position.  Its geometric
significance will be discussed in detail later.

We have now established the relationship between path
independence and the existence of a scalar function
$\psi(x,\,y,\,z)$ such that $\bfF = \nabla\psi$.  Since there is also
a relationship between path independence and the fact that
$\nabla\times\bfF = 0$, you may suspect that $\nabla\times\bfF = 0$
and $\bfF = \nabla\psi$ are also related.  Indeed, if
$\bfF = \nabla\psi$, then under suitable conditions,
$\nabla\times\bfF = 0$.  This is easily established.  Consider, for
example, the $x$-component of $\nabla\times\bfF$:
\[
(\nabla\times\bfF)_x
= \pd{F_z}{y} - \pd{F_y}{z}
= \frac{\partial}{\partial y}\!\left(\pd{\psi}{z}\right)
- \frac{\partial}{\partial z}\!\left(\pd{\psi}{y}\right)
= \pdd{\psi}{y\,\partial z} - \pdd{\psi}{z\,\partial y} = 0.
\]

This last equality follows if $\psi$ and its first and second
derivatives are continuous, for then
$\partial^2\psi/\partial y\,\partial z = \partial^2\psi/\partial z\,\partial y$.
Obviously, the other two components of $\nabla\times\bfF$ can be
shown to vanish in exactly the same way.  Thus,
\[
F_q = \pd{\psi}{q}\qquad (q = x,\,y,\,z)
\qquad\Rightarrow\qquad
\nabla\times\bfF = 0.
\]

The converse of what we have just shown would assert that if
$\nabla\times\bfF = 0$, then there exists a scalar function $\psi$
such that $\bfF = \nabla\psi$, a statement that is true, provided the
region of interest is simply connected.  To understand this, we can
consult Figure~IV--3, which shows how path independence of the
line integral of $\bfF \cdot \hat{\mathbf{t}}$,
$\nabla\times\bfF = 0$, and $\bfF = \nabla\psi$ are related.  The
solid arrows in the diagram represent implications that hold in
general, provided $\bfF$ is smooth.  The dashed arrows represent
implications requiring not only that $\bfF$ be smooth, but that the
region of interest be simply connected.  We have already shown that
(1) implies both (2) and (3) and that (3) implies (1) in a simply
connected region.  Combining these two statements, we see that (3)
implies (2) in a simply connected region.

\begin{figure}[ht]
\centering
\includegraphics[width=0.6\textwidth]{figures/ch04/fig_IV_3.png}
\caption{}
\label{fig:IV-3}
\end{figure}

In practice, just as the functions we deal with usually have
continuous first derivatives (and are therefore smooth), the regions
we work with are simply connected.  In such circumstances we can
relax a bit and regard the three statements summarized in Figure~IV--3 as equivalent: each implies and is implied by each of
the others.  However, you should be aware of simple connectedness
and its implications for the relations among the three statements.

To give a simple example of the ideas we have been discussing,
consider the vector function
\[
\bfF(x,\,y,\,z) = \bfi\,y + \bfj\,x.
\]
This function is smooth everywhere, and we have already noted that
its curl is zero (page~91).  According to what we have just said,
this means there must be a scalar function $\psi(x,\,y,\,z)$ such
that $\bfF$ is its gradient.  Thus, $\psi$ must satisfy
\[
F_x = y = \pd{\psi}{x}\,,\qquad
F_y = x = \pd{\psi}{y}\,,\qquad
F_z = 0 = \pd{\psi}{z}\,.
\]
Clearly, $\psi(x,\,y,\,z) = xy + C$, where $C$ is an arbitrary
constant, satisfies these relations.  This should be contrasted with
the case of the function $\bfF = \bfi\,y - \bfj\,x$, the curl of which
does \emph{not} vanish (page~92).  If this function were the
gradient of a scalar function $\psi$, we should have
\[
F_x = y = \pd{\psi}{x}\,,\qquad
F_y = -x = \pd{\psi}{y}\,,\qquad
F_z = 0 = \pd{\psi}{z}\,,
\]
but, as you should be able to convince yourself, there is no function
$\psi$ that satisfies these three equations.

The expression we have written for the gradient of a scalar
function $\psi(x,\,y,\,z)$, namely,
\[
\nabla\psi = \bfi\,\pd{\psi}{x} + \bfj\,\pd{\psi}{y}
            + \bfk\,\pd{\psi}{z}\,,
\]
is really just the form of this operator in Cartesian coordinates.
To find the form of the gradient in other coordinate systems, if you
go about it straightforwardly, is a tedious job.  For example, to
find the gradient in cylindrical coordinates, we would first have to
express the Cartesian unit vectors $\bfi$, $\bfj$, and $\bfk$ in terms
of the analogous quantities $\hat{\mathbf{e}}_r$,
$\hat{\mathbf{e}}_\theta$, and $\hat{\mathbf{e}}_z$ in cylindrical
coordinates.  Then, using $x = r\cos\theta$, $y = r\sin\theta$, and
the chain rule for differentiation, we would have to express
derivatives with respect to $x$, $y$, and $z$ in terms of those with
respect to $r$, $\theta$, and $z$.  We shall not pursue this matter
here because later (see pages~141~ff.) an easier and faster method
will be available to us.  For the present we merely quote the form of
the gradient in cylindrical and in spherical coordinates.

\emph{Cylindrical:}
\begin{equation}
\nabla\psi
= \hat{\mathbf{e}}_r\,\pd{\psi}{r}
+ \hat{\mathbf{e}}_\theta\,\frac{1}{r}\,\pd{\psi}{\theta}
+ \hat{\mathbf{e}}_z\,\pd{\psi}{z}.
\label{eq:IV-3}
\end{equation}

\emph{Spherical:}
\begin{equation}
\nabla\psi
= \hat{\mathbf{e}}_r\,\pd{\psi}{r}
+ \hat{\mathbf{e}}_\phi\,\frac{1}{r}\,\pd{\psi}{\phi}
+ \hat{\mathbf{e}}_\theta\,\frac{1}{r\sin\phi}\,\pd{\psi}{\theta}
\label{eq:IV-4}
\end{equation}

A coordinate-free definition of the gradient analogous to the ones
given for the divergence [Equation~(II--17)] and the curl
[Equation~(III--8)] is discussed in Problem~IV--25.


%% ================================================================
\section{Finding the Electrostatic Field}
%% ================================================================

We began our discussion of vector calculus with a search for some
convenient method for finding the electrostatic field.  Our
investigations led us to the differential form of Gauss' law,
\[
\nabla\cdot\bfE = \rho/\epszero.
\]

Even this expression is not often useful for finding $\bfE$ because
it is one equation in three unknowns ($E_x$, $E_y$, and $E_z$ in
Cartesian coordinates).  Now, at last, we are able to complete our
discussion and write down the equations that are often the most
useful of all known methods for finding the field.

This final step rests on the observation that since
\[
\oint_C \bfE \cdot \hat{\mathbf{t}}\ds = 0
\]
for any closed path $C$, the field $\bfE$ can be written as the
gradient of a scalar function.  Conventionally, this function, called
the electrostatic potential, is designated $\Phi(x,\,y,\,z)$, and we
write\footnote{The negative sign in this equation is not put there just
to make life more difficult; there is a good reason for it.  See the
discussion on page~139.}
\[
\bfE = -\nabla\Phi.
\]

Combining this equation with the differential form of Gauss' law
[Equation~(II--17)], we get
\[
\nabla\cdot(-\nabla\Phi) = \rho/\epszero,
\]
or
\[
\nabla\cdot(\nabla\Phi) = -\rho/\epszero.
\]

When we write out the left-hand side of this equation in detail, we
find
\[
\nabla\cdot(\nabla\Phi)
= \left(\bfi\,\frac{\partial}{\partial x}
  + \bfj\,\frac{\partial}{\partial y}
  + \bfk\,\frac{\partial}{\partial z}\right)
  \cdot\left(\bfi\,\pd{\Phi}{x}
  + \bfj\,\pd{\Phi}{y}
  + \bfk\,\pd{\Phi}{z}\right)
= \pdd{\Phi}{x} + \pdd{\Phi}{y} + \pdd{\Phi}{z}\,,
\]
and so
\begin{equation}
\pdd{\Phi}{x} + \pdd{\Phi}{y} + \pdd{\Phi}{z}
= -\rho/\epszero.
\label{eq:IV-5}
\end{equation}

Equation~(\ref{eq:IV-5}) can be written more compactly by
introducing a new operator, called the \emph{Laplacian}, which is
denoted, for fairly obvious reasons, by the symbol $\nabla^2$ (read
``del squared'').  That is,
\begin{equation}
\nabla^2 = \nabla\cdot\nabla
= \left(\bfi\,\frac{\partial}{\partial x}
  + \bfj\,\frac{\partial}{\partial y}
  + \bfk\,\frac{\partial}{\partial z}\right)
  \cdot\left(\bfi\,\frac{\partial}{\partial x}
  + \bfj\,\frac{\partial}{\partial y}
  + \bfk\,\frac{\partial}{\partial z}\right)
= \pdd{}{x} + \pdd{}{y} + \pdd{}{z}\,.
\label{eq:IV-6}
\end{equation}

In this new notation Equation~(\ref{eq:IV-5}) becomes
\begin{equation}
\laplacian\Phi = -\rho/\epszero.
\label{eq:IV-7}
\end{equation}

Equation~(\ref{eq:IV-6}) provides the form of the Laplacian in
Cartesian coordinates; its forms in cylindrical and spherical
coordinates will be given in the next section.  The best definition of
the Laplacian is probably
\[
\laplacian f = \nabla\cdot(\nabla f),
\]
where $f$ is some suitably continuous scalar function of position.
This definition has the important advantage of being independent of
the coordinate system.

Equation~(\ref{eq:IV-7}) is called Poisson's equation.  It is a
linear, second-order partial differential equation in \emph{one}
unknown, the scalar function $\Phi(x,\,y,\,z)$, and is the
culmination of our long search for a method of determining the
electrostatic field.  A great body of work exists describing the many
elegant mathematical schemes that have been devised to solve it,
and a few simple examples are given in the next section.  In any
problem, once we have $\Phi$, the field is trivial to find using
$\bfE = -\nabla\Phi$.

At any point in space where there is no electric charge, the
density $\rho$ is zero and Poisson's equation reduces to
\[
\laplacian\Phi = 0.
\]
This is called \emph{Laplace's equation} and is more often used than
Poisson's equation.  The reason for this is that usually charges are
distributed over various objects; this gives rise to a field, and we
are interested in finding the potential (and from it, the field) in the
charge-free space between the objects.  In the simplest of situations
it is possible to specify ``boundary conditions,'' that is, the

\begin{figure}[ht]
\centering
\includegraphics[width=0.6\textwidth]{figures/ch04/fig_IV_4.png}
\caption{}
\label{fig:IV-4}
\end{figure}

\noindent value of the potential on the surfaces of these objects
(Figure~IV--4).  We then find that solution of Laplace's equation
which takes on the given values on the surfaces.  This is
illustrated in the next section.


%% ================================================================
\section{Using Laplace's Equation}
%% ================================================================

Whether solving Laplace's equation is or is not a topic in vector
calculus is a moot point, but the basis of our entire discussion has
been a search for a method to calculate electric fields.  Since
Laplace's equation is the end product of that search, we can
scarcely omit a few examples to show how it works.\footnote{This
section is not essential to what follows and may be omitted.}

We begin with an especially simple problem.  Imagine we have two
very large (``infinite'') parallel plates separated by a distance $s$
(Figure~IV--5).  Choosing a coordinate system as shown in the
figure, let the plate at $x = 0$ be held at zero potential and that at
$x = s$ at $V_0$.  Our object is to find the potential and the
electric field in the space between the two plates.  Because the
plates are infinitely large, there is nothing to distinguish a point
$(x,\,y,\,z)$ from any other point $(x,\,y',\,z')$ having the same
$x$-coordinate.  It follows that the potential $\Phi$ depends on $x$
but not on $y$ or $z$.  Thus, $\laplacian\Phi$ reduces in this case
simply to $d^2\Phi/dx^2$, and so Laplace's equation and the
associated boundary conditions are

\begin{figure}[ht]
\centering
\includegraphics[width=0.6\textwidth]{figures/ch04/fig_IV_5.png}
\caption{}
\label{fig:IV-5}
\end{figure}

\[
\frac{d^2\Phi}{dx^2} = 0
\]
and
\[
\Phi = \begin{cases} 0 & \text{at } x = 0 \\ V_0 & \text{at } x = s. \end{cases}
\]

This is a trivial problem and the solution is
\[
\Phi(x) = \frac{V_0 x}{s}\,.
\]

The electric field is found using $\bfE = -\nabla\Phi$, which yields
\[
E_x = \frac{-V_0}{s}\qquad\text{and}\qquad E_y = E_z = 0.
\]

Thus, the field is a constant vector normal to the plates.  This is an
excellent approximation to the potential and field between, but far
from the edges of, two plates whose linear dimensions are large
compared with their separation.  You may recognize this
arrangement as a parallel plate capacitor.

Our second example is a spherical capacitor, that is, two
concentric spheres having radii $R_1$ and $R_2$ with the inner one
maintained at a potential $V_0$ and the outer at zero
(Figure~IV--6).  We are required to find the potential and field
everywhere between

\begin{figure}[ht]
\centering
\includegraphics[width=0.6\textwidth]{figures/ch04/fig_IV_6.png}
\caption{}
\label{fig:IV-6}
\end{figure}

\noindent the spheres.  In this situation we would obviously do well
to work in spherical coordinates $r$, $\theta$, and $\phi$, in which
Laplace's equation between the spheres has the imposing form
\[
\laplacian\Phi
= \frac{1}{r^2}\frac{\partial}{\partial r}
  \!\left(r^2\pd{\Phi}{r}\right)
+ \frac{1}{r^2\sin\phi}\frac{\partial}{\partial\phi}
  \!\left(\sin\phi\,\pd{\Phi}{\phi}\right)
+ \frac{1}{r^2\sin^2\!\phi}\,\pdd{\Phi}{\theta} = 0.
\]
(See Problem~IV--23.)  Fortunately, we need not work with this
equation as it stands; a little thought will convince you that $\Phi$
can only be a function of $r$, since there is no way to distinguish a
point $(r,\,\theta,\,\phi)$ from another $(r,\,\theta',\,\phi')$ with the
same $r$ but different $\theta$ and $\phi$.  Thus,
\[
\pd{\Phi}{\theta} = \pd{\Phi}{\phi} = 0,
\]
and Laplace's equation reduces to
\begin{equation}
\frac{1}{r^2}\frac{d}{dr}\!\left(r^2\frac{d\Phi}{dr}\right) = 0.
\label{eq:IV-8}
\end{equation}

We are interested in the solution of this equation that is valid for
$R_1 < r < R_2$ and satisfies the boundary conditions
\[
\Phi(r) = \begin{cases} V_0 & \text{at } r = R_1 \\ 0 & \text{at } r = R_2. \end{cases}
\]

Multiplying Equation~(\ref{eq:IV-8}) by $r^2$ and putting
$\psi = d\Phi/dr$, we get
\[
\frac{d}{dr}(r^2\psi) = 0,
\]
and so
\[
r^2\psi = c_1,
\]
where $c_1$ is a constant.  Hence,
\[
\psi = \frac{d\Phi}{dr} = \frac{c_1}{r^2}\,,
\]
and it follows that
\begin{equation}
\Phi = -\frac{c_1}{r} + c_2,
\label{eq:IV-9}
\end{equation}
where $c_2$ is another constant.  Imposing the boundary conditions,
we find
\[
-\frac{c_1}{R_1} + c_2 = V_0
\qquad\text{and}\qquad
-\frac{c_1}{R_2} + c_2 = 0,
\]
whence
\[
c_1 = \frac{V_0 R_1 R_2}{R_1 - R_2}
\qquad\text{and}\qquad
c_2 = \frac{V_0 R_1}{R_1 - R_2}\,.
\]

Substituting these in the expression for the potential
[Equation~(\ref{eq:IV-9})], we get
\[
\Phi(r) = \frac{V_0 R_1}{R_1 - R_2}
  \left(1 - \frac{R_2}{r}\right),
\qquad R_1 < r < R_2.
\]

To get the electric field, we must take the gradient of $\Phi$, and
this is clearly most conveniently done in spherical coordinates [see
Equation~(\ref{eq:IV-4})].  However, since in this case $\Phi$
depends only on $r$, we get only a radial component:
\[
E_r = -\frac{d\Phi}{dr}
    = -\frac{V_0 R_1 R_2}{R_1 - R_2}\,\frac{1}{r^2}\,,
\]
\[
E_\theta = E_\phi = 0,\qquad (R_1 < r < R_2).
\]

Our third and last example is more complicated (and more
interesting) than the foregoing.  If a potential difference is
maintained between two ``infinite'' parallel plates $P$ and $P'$
(Figure~IV--7), then we know from our first example that the field
between them is a constant vector normal to the plates.  Choosing a
coordinate system as shown in the figure (with the $z$-axis out of
the plane of the paper), we have $\bfE = E_0\bfi$, where $E_0$ is a
constant.  Let an ``infinitely'' long cylinder held at zero potential
be situated between the plates with its axis along the $z$-axis.  Let
its radius $R$ be small compared with the plate separation.  What
are the potential and the electric field outside the cylinder and
between the plates?  Here, clearly, we should use cylindrical
coordinates $(r,\,\theta,\,z)$, in which case Laplace's equation reads

\begin{figure}[ht]
\centering
\includegraphics[width=0.6\textwidth]{figures/ch04/fig_IV_7.png}
\caption{}
\label{fig:IV-7}
\end{figure}

\[
\laplacian\Phi
= \frac{1}{r}\frac{\partial}{\partial r}
  \!\left(r\,\pd{\Phi}{r}\right)
+ \frac{1}{r^2}\,\pdd{\Phi}{\theta}
+ \pdd{\Phi}{z} = 0.
\]
(See Problem~IV--21.)  You should convince yourself that $\Phi$ in
this case must be independent of $z$, so this equation simplifies
somewhat to
\begin{equation}
\frac{1}{r}\frac{\partial}{\partial r}
  \!\left(r\,\pd{\Phi}{r}\right)
+ \frac{1}{r^2}\,\pdd{\Phi}{\theta} = 0.
\label{eq:IV-10}
\end{equation}

There are two boundary conditions, of which the first is
\[
\Phi(r,\,\theta) = 0\qquad\text{at}\quad r = R.
\]

The second condition has to do with the fact that at large values of
$r$, the influence of the cylinder is negligible and the field must
be, to a good approximation, what it would be if the cylinder were
not present at all, that is, $E_0\bfi$.  To put this in terms of the
potential, we note that
\[
\Phi = -E_0 x
\]
will provide just such a field.  Since $x = r\cos\theta$, we can
write the second boundary condition
\begin{equation}
\Phi(r,\,\theta) = -E_0 r\cos\theta,\qquad r \gg R.
\label{eq:IV-11}
\end{equation}

Let's try to solve Laplace's equation for this problem
[Equation~(\ref{eq:IV-10})] by assuming we can write
\begin{equation}
\Phi(r,\,\theta) = f(r)\cos\theta,
\label{eq:IV-12}
\end{equation}
where $f(r)$ is an as yet unknown function.  What prompts us to do
this is the fact that the second boundary condition
[Equation~(\ref{eq:IV-11})] has precisely this form---a function of
$r$ multiplied by $\cos\theta$.  If we substitute
Equation~(\ref{eq:IV-12}) into Equation~(\ref{eq:IV-10}), the
result is a differential equation for the function $f(r)$:
\[
\frac{d^2 f}{dr^2} + \frac{1}{r}\frac{df}{dr} - \frac{1}{r^2}\,f = 0.
\]

Putting $f(r) = r^\lambda$ where $\lambda$ is a constant leads to
\[
\lambda(\lambda - 1)r^{\lambda-2} + \lambda r^{\lambda-2} - r^{\lambda-2} = 0,
\]
or
\[
\lambda^2 = 1,
\]
and $\lambda = \pm 1$.  Hence we get
\[
f(r) = Ar + \frac{B}{r}\,,
\]
where $A$ and $B$ are constants.  Thus, our solution is
\[
\Phi(r,\,\theta) = \left(Ar + \frac{B}{r}\right)\cos\theta.
\]

The first boundary condition requires that
\[
AR + \frac{B}{R} = 0,
\]
or
\[
B = -AR^2.
\]

Hence,
\[
\Phi(r,\,\theta) = Ar\cos\theta - \frac{AR^2}{r}\cos\theta.
\]

To impose the second condition, we note that for $r$ large, the
second term in this last equation is negligible compared with the
first.  Thus,
\[
\Phi(r,\,\theta) \simeq Ar\cos\theta,\qquad r\text{ large.}
\]

We satisfy the second boundary condition by choosing $A = -E_0$.
The complete solution is thus
\[
\Phi(r,\,\theta) = -E_0 r\left(1 - \frac{R^2}{r^2}\right)\cos\theta.
\]

To find the electric field, we proceed as usual with
$\bfE = -\nabla\Phi$.  Using Equation~(\ref{eq:IV-3}), we get
\[
E_r = -\pd{\Phi}{r}
    = E_0\left[1 + \left(\frac{R}{r}\right)^{\!2}\right]\cos\theta,
\]
\[
E_\theta = -\frac{1}{r}\pd{\Phi}{\theta}
         = -E_0\left[1 - \left(\frac{R}{r}\right)^{\!2}\right]\sin\theta,
\]
\[
E_z = -\pd{\Phi}{z} = 0.
\]

You should verify that for large $r$, this field reduces to
$E_0\bfi$ as required.

You may find this last example disquieting, since a certain amount
of clever guesswork is used in finding the potential.  Actually,
there are standard procedures that, in problems of this kind, lead
more or less straightforwardly to the solution.  A discussion of
these procedures, however, would be very lengthy and (in the
well-worn phrase) beyond the scope of this text.  Before moving on,
however, one further point is worth making: A solution of Laplace's
equation that satisfies appropriate boundary conditions is
\emph{unique}.  That is to say, there is one and only one such
solution, so that if we solve a problem by guesswork and
skullduggery, and someone else solves it with refined and elegant
mathematical techniques, the two solutions, in spite of their
disparate pedigrees, must be the same.  In Problem~IV--24 you will
be led through a proof of this remarkable fact.


%% ================================================================
\section{Directional Derivatives and the Gradient}
%% ================================================================

We have introduced the gradient as a sort of mathematical artifice
useful in discussing path-independent line integrals.  We now turn
to a more detailed examination of the gradient in order to describe
its geometrical significance.

Before beginning our discussion, we make a few comments on
Taylor series, since these are needed in what follows.  For a scalar
function of one variable that is suitably continuous and
differentiable, we have
\[
f(x + \Delta x) = f(x) + \Delta x\,f'(x)
+ \tfrac{1}{2}(\Delta x)^2 f''(x) + \cdots\,.
\]

This says that the value of the function at some point
$x + \Delta x$ can be written as the sum of (usually) infinitely many
terms that involve the function and its derivatives at some other
point $x$.  Among other things, this Taylor series is useful for
calculation, for if the two points are close together (that is, if
$\Delta x$ is small), then we can truncate the series after a certain
number of terms (which we hope is small), since the neglected
terms, each proportional to some large power of the small number
$\Delta x$, will sum to a value that is negligible.

Taylor series can also be formed for functions of several
variables.  Thus, for a function of two variables we have
\begin{equation}
f(x + \Delta x,\,y + \Delta y)
= f(x,\,y) + \Delta x\,\pd{f}{x} + \Delta y\,\pd{f}{y} + \cdots\,.
\label{eq:IV-13}
\end{equation}

This says that the value of the function at some point
$(x + \Delta x,\,y + \Delta y)$ can be written as a sum of (usually)
infinitely many terms that involve the function and its derivatives at
some other point $(x,\,y)$.  We shall never need the explicit form of
the remaining terms of this series [represented by the dots in
Equation~(\ref{eq:IV-13})].  We should know, however, that these
terms involve higher powers of the ``small'' numbers $\Delta x$ and
$\Delta y$ (for example, $\Delta x^2$, $\Delta y^2$,
$\Delta x\,\Delta y$, $\Delta x^3$, $\Delta y^3$,
$\Delta x^2\,\Delta y$, and so on).  With these simple ideas in mind
we turn now to our main task.

Consider some function $z = f(x,\,y)$.  Geometrically, this
represents a surface as shown in Figure~IV--8(a).  Let $(x,\,y)$ be
the coordinates of a point $P$ in the $xy$-plane.  The height of
the surface above this point is represented by the length of the
dotted line $PQ$; that is, $PQ = z = f(x,\,y)$.  Suppose now we
take a short step in the $xy$-plane to a new point $P'$ with
coordinates $(x + \Delta x,\,y + \Delta y)$.  The height of the
surface above this point is
$P'Q' = f(x + \Delta x,\,y + \Delta y)$.  Let $\Delta s$ be the
length of the step ($\Delta s = PP'$).

\begin{figure}[ht]
\centering
\includegraphics[width=0.6\textwidth]{figures/ch04/fig_IV_8a.png}
\caption{}
\label{fig:IV-8a}
\end{figure}

We next ask how much the function $f$ has changed as a result of
taking this step.  Clearly, this change is the difference in the two
heights $PQ$ and $P'Q'$, and
\[
P'Q' - PQ \equiv \Delta f
= f(x + \Delta x,\,y + \Delta y) - f(x,\,y).
\]

Applying the Taylor series formula stated above
[Equation~(\ref{eq:IV-13})], we get
\[
\Delta f = f(x,\,y) + \Delta x\,\pd{f}{x} + \Delta y\,\pd{f}{y}
+ \cdots - f(x,\,y)
= \Delta x\,\pd{f}{x} + \Delta y\,\pd{f}{y} + \cdots\,.
\]

We now recast this expression by what at first may seem an
unnecessary elaboration of the notation.  Let $\Delta\mathbf{s}$ be a
vector that has magnitude $\Delta s$ and points from $P$ to $P'$.
Clearly,
\[
\Delta\mathbf{s} = \bfi\,\Delta x + \bfj\,\Delta y.
\]

But the gradient of $f$ is
\[
\nabla f = \bfi\,\pd{f}{x} + \bfj\,\pd{f}{y}
\]
(an obvious specialization of the gradient notation to a function of
two, rather than three, variables).  It follows at once that
\[
\Delta f = (\Delta\mathbf{s})\cdot(\nabla f) + \cdots\,.
\]

Complicating matters slightly more, let $\hat{\mathbf{u}}$ be a unit
vector in the direction of $\Delta\mathbf{s}$.  Then
\[
\Delta\mathbf{s} = \hat{\mathbf{u}}\,\Delta s
\]
and
\[
\Delta f = (\hat{\mathbf{u}}\cdot\nabla f)\,\Delta s + \cdots\,,
\]
so that
\[
\frac{\Delta f}{\Delta s} = \hat{\mathbf{u}}\cdot\nabla f + \cdots\,.
\]

We now take the limit of this equation to get
\begin{equation}
\frac{df}{ds}
\equiv \lim_{\Delta s\to 0}\frac{\Delta f}{\Delta s}
= \hat{\mathbf{u}}\cdot\nabla f.
\label{eq:IV-14}
\end{equation}

There is no longer any need for ``$+\,\cdots$,'' since the dots
represented terms that go to zero as $\Delta s$ goes to zero.

This new expression [Equation~(\ref{eq:IV-14})] has a simple
interpretation: it is the \emph{rate} of change of the function
$f(x,\,y)$ in the direction of $\Delta\mathbf{s}$ (that is, of
$\hat{\mathbf{u}}$).  Redrawing Figure~IV--8(a) and passing a plane
through $P$ and $P'$ parallel to the $z$-axis
[Figure~IV--8(b)], we see that it cuts the surface $z = f(x,\,y)$ in
a curve $C$.  The quantity $df/ds$ defined in
Equation~(\ref{eq:IV-14}) is the slope of this curve at the point
$Q$.

\begin{figure}[ht]
\centering
\includegraphics[width=0.6\textwidth]{figures/ch04/fig_IV_8b.png}
\caption{}
\label{fig:IV-8b}
\end{figure}

The quantity $df/ds$ is called the \emph{directional derivative} of
$f$.  Although the analysis given earlier that led to this derivative
was for functions of two variables, the results all apply to functions
of three (or more) variables.  Thus,
\[
\frac{d}{ds}\,F(x,\,y,\,z) = \hat{\mathbf{u}}\cdot\nabla F
\]
is the rate of change of the function $F(x,\,y,\,z)$ in the direction
specified by the unit vector $\hat{\mathbf{u}}$.

An example of the directional derivative may be amusing here.
We'll work with a function of two variables so that we can draw
pictures.  Thus, let's consider
\[
z = f(x,\,y) = (x^2 + y^2)^{1/2},
\]
which is an inverted right circular cone whose axis coincides with
the $z$-axis [see Figure~IV--9(a)].  We ask for the directional
derivative of this function at some point $x = a$ and $y = b$ and in
the direction specified by
$\hat{\mathbf{u}} = \bfi\cos\theta + \bfj\sin\theta$ [see
Figure~IV--9(b)].  First we need the gradient of $f(x,\,y)$.  But
\[
\pd{f}{x} = \frac{x}{z}\qquad\text{and}\qquad
\pd{f}{y} = \frac{y}{z}\,,
\]
as you can easily verify.  Thus,
\[
\nabla f = \frac{\bfi\,x + \bfj\,y}{z}
\]

\begin{figure}[ht]
\centering
\includegraphics[width=0.45\textwidth]{figures/ch04/fig_IV_9a.png}
\caption{}
\label{fig:IV-9a}
\end{figure}

\begin{figure}[ht]
\centering
\includegraphics[width=0.45\textwidth]{figures/ch04/fig_IV_9b.png}
\caption{}
\label{fig:IV-9b}
\end{figure}

\noindent and
\[
\frac{df}{ds} = \hat{\mathbf{u}}\cdot\nabla f
= \frac{x\cos\theta + y\sin\theta}{z}
\xrightarrow{\text{at }(a,b)}
\frac{a\cos\theta + b\sin\theta}{\sqrt{a^2+b^2}}\,.
\]

Suppose $\theta$ is chosen so that $\hat{\mathbf{u}}$ is in the radial
direction as indicated in Figure~IV--9(c).  This means

\begin{figure}[ht]
\centering
\includegraphics[width=0.45\textwidth]{figures/ch04/fig_IV_9c.png}
\caption{}
\label{fig:IV-9c}
\end{figure}

\[
\cos\theta = \frac{a}{(a^2+b^2)^{1/2}}\,,\qquad
\sin\theta = \frac{b}{(a^2+b^2)^{1/2}}\,,
\]
and so
\[
\frac{df}{ds}
= \frac{a}{\sqrt{a^2+b^2}}\cdot\frac{a}{\sqrt{a^2+b^2}}
+ \frac{b}{\sqrt{a^2+b^2}}\cdot\frac{b}{\sqrt{a^2+b^2}}
= 1.
\]

\begin{figure}[ht]
\centering
\includegraphics[width=0.55\textwidth]{figures/ch04/fig_IV_9d.png}
\caption{}
\label{fig:IV-9d}
\end{figure}

The significance of this result is brought out in Figure~IV--9(d).

A second interesting case is that in which $\hat{\mathbf{u}}$ is chosen
perpendicular to the direction of the previous example [see
Figure~IV--9(e)].  We then have

\begin{figure}[ht]
\centering
\includegraphics[width=0.45\textwidth]{figures/ch04/fig_IV_9e.png}
\caption{}
\label{fig:IV-9e}
\end{figure}

\[
\cos\theta = \frac{-b}{(a^2+b^2)^{1/2}}\,,\qquad
\sin\theta = \frac{a}{(a^2+b^2)^{1/2}}\,,
\]
and so
\[
\frac{df}{ds}
= \frac{a}{\sqrt{a^2+b^2}}\left(\frac{-b}{\sqrt{a^2+b^2}}\right)
+ \frac{b}{\sqrt{a^2+b^2}}\left(\frac{a}{\sqrt{a^2+b^2}}\right)
= 0.
\]

\begin{figure}[ht]
\centering
\includegraphics[width=0.55\textwidth]{figures/ch04/fig_IV_9f.png}
\caption{}
\label{fig:IV-9f}
\end{figure}

The meaning of this result is illustrated in Figure~IV--9(f).


%% ================================================================
\section{Geometric Significance of the Gradient}
%% ================================================================

With the concept of the directional derivative at our disposal, we
are now in a position to give a geometric interpretation of the
gradient.  At some point $P_0$ with coordinates $(x_0,\,y_0,\,z_0)$
we have
\[
\left(\frac{dF}{ds}\right)_{\!0} = \hat{\mathbf{u}}\cdot(\nabla F)_0,
\]
where the subscript $0$ means the quantity is to be evaluated at the
point $(x_0,\,y_0,\,z_0)$.  Now $(\nabla F)_0$, the gradient of $F$
evaluated at $P_0$, may be represented by an arrow emanating from
that point as shown in Figure~IV--10(a).  If we ask in what
direction we must move to make $(dF/ds)_0$ as large as possible, it
is clear that $\hat{\mathbf{u}}$ should be in the same direction as
$(\nabla F)_0$.  This is because if we let $\alpha$ be the angle
between $\hat{\mathbf{u}}$ and $(\nabla F)_0$, then
$(dF/ds)_0 = |\nabla F|_0\cos\alpha$, and this is as large as it can
be when $\alpha = 0$.  Thus, \emph{the gradient of a scalar function}
$F(x,\,y,\,z)$ \emph{is a vector that is in the direction in which}
$F$ \emph{undergoes the greatest rate of increase and that has
magnitude equal to the rate of increase in that direction.}

\begin{figure}[ht]
\centering
\includegraphics[width=0.6\textwidth]{figures/ch04/fig_IV_10a.png}
\caption{}
\label{fig:IV-10a}
\end{figure}

To illustrate this interpretation of the gradient, let us go back to
the inverted cone $z = f(x,\,y) = (x^2+y^2)^{1/2}$ we discussed
earlier.  We learned that
\[
\nabla f = \frac{\bfi\,x + \bfj\,y}{z}
\]
and
\[
\frac{df}{ds}
= \frac{a\cos\theta + b\sin\theta}{\sqrt{a^2+b^2}} \equiv D(\theta).
\]

To find the direction in which $f(x,\,y)$ undergoes the greatest
rate of change, we set
\[
\frac{dD}{d\theta} = \frac{-a\sin\theta + b\cos\theta}{\sqrt{a^2+b^2}} = 0.
\]

This gives $\tan\theta = b/a$, whence
$\cos\theta = a/(a^2+b^2)^{1/2}$ and
$\sin\theta = b/(a^2+b^2)^{1/2}$.  So
$(df/ds)_{\max} = 1$.  On the other hand,
\[
|\nabla f|
= \left[\frac{x^2+y^2}{z^2}\right]^{1/2} = 1,
\]
since $z^2 = x^2 + y^2$.  Furthermore, $\tan\theta = b/a$
corresponds to the direction $a\bfi + b\bfj$, while at the point
$(a,\,b)$,
\[
\nabla f = \frac{a\bfi + b\bfj}{(a^2+b^2)^{1/2}}\,,
\]
which is a vector in the same direction.  Thus both properties of the
gradient are illustrated; it's in the direction of maximum rate of
increase, and its magnitude is equal to the rate of increase in that
direction.

As a second example of the interpretation of the gradient, we
consider the plane $z = f(x,\,y) = 1 - x - y$ shown in
Figure~IV--10(b).  It's easy to see that
$\nabla f = -\bfi - \bfj$.  Using
$\hat{\mathbf{u}} = \bfi\cos\theta + \bfj\sin\theta$ as before, we
find $df/ds = \hat{\mathbf{u}}\cdot\nabla f = -\cos\theta - \sin\theta \equiv D(\theta)$.
Thus
\[
\frac{dD}{d\theta} = \sin\theta - \cos\theta = 0,
\]

\begin{figure}[ht]
\centering
\includegraphics[width=0.55\textwidth]{figures/ch04/fig_IV_10b.png}
\caption{}
\label{fig:IV-10b}
\end{figure}

\noindent which yields $\theta = \pi/4$ or $5\pi/4$.  The second
derivative test shows that $\pi/4$ corresponds to a minimum and
$5\pi/4$, to a maximum.  It can be seen from Figure~IV--10(b) that
the greatest rate of increase is indeed at an angle of $5\pi/4$.
Moreover, the greatest rate of increase is
\[
\left(\frac{df}{ds}\right)_{\!\max} = D(5\pi/4) = \sqrt{2}\,,
\]
whereas $|\nabla f| = |-\bfi - \bfj| = \sqrt{2}$.  Again, both
properties of the gradient are illustrated by this example.

With this geometric interpretation of the gradient at our disposal,
we can now see the reason for the negative sign in the equation
$\bfE = -\nabla\Phi$: Since $\nabla\Phi$ is a vector in the direction
of increasing $\Phi$, the force on a \emph{positive} charge $q$ is
$\bfF = q\bfE = -q\nabla\Phi$, which is in the direction of
\emph{decreasing} $\Phi$.  Thus, the negative sign ensures that a
positive charge moves ``downhill'' from a higher to a lower
potential.

There is another property of the gradient useful in understanding
its geometric significance.  To make this discussion concrete, let
$T(x,\,y,\,z)$ be a scalar function that gives the temperature at

\begin{figure}[ht]
\centering
\includegraphics[width=0.6\textwidth]{figures/ch04/fig_IV_11.png}
\caption{}
\label{fig:IV-11}
\end{figure}

\noindent any point $(x,\,y,\,z)$.  The locus of all points having the
same temperature $T_0$ is (in the simplest case) a surface whose
equation is $T(x,\,y,\,z) = T_0$ (Figure~IV--11).  This is called an
isothermal surface.  We now show that $\nabla T$ is a vector normal
to the isothermal surface.  Let $C$ be any curve lying in the
isothermal surface and let $P$ be any point on $C$.  Let
$\hat{\mathbf{u}}$ be the unit vector tangent to $C$ at $P$ (it doesn't
matter which direction along $C$ we take).  The directional
derivative in the direction $\hat{\mathbf{u}}$ is
\[
\left(\frac{dT}{ds}\right) = \hat{\mathbf{u}}\cdot\nabla T = 0
\]
because $T$ does not change as we move along the isothermal
surface.  If the scalar product of two vectors, neither of them zero,
vanishes, the two vectors are perpendicular.  Thus $\nabla T$ is
perpendicular to $C$ at $P$.  By the same argument it is
perpendicular to \emph{any} curve on the surface through $P$ (such
as $C'$ in Figure~IV--11).  But this can be true only if $\nabla T$
is normal to the isothermal surface at $P$.  In general then,
$\nabla f(x,\,y,\,z)$, \emph{where} $f(x,\,y,\,z)$ \emph{is a scalar
function, is normal to the surface}
$f(x,\,y,\,z) = \textit{constant}$.\footnote{The connection between
this property of the gradient and our earlier expression for the unit
vector normal to a surface [Equation~(II--4)] is the subject of
Problem~IV--20.}

A simple example of this property of the gradient is provided by
the function $F(x,\,y,\,z) = x^2 + y^2 + z^2$.  The surface
$F(x,y,z) = \text{constant}$ is, of course, a sphere (assuming the
constant is positive).  As you should verify for yourself,
$\nabla F = 2(\bfi\,x + \bfj\,y + \bfk\,z) = 2\bfr$.  Thus, we have
a familiar result: A vector normal to a spherical surface is in the
radial direction.  We'll leave it to you to ponder the geometric
relation between the electrostatic field $\bfE$ and its equipotential
surfaces $\Phi(x,\,y,\,z) = \text{constant}$.

\begin{figure}[ht]
\centering
\includegraphics[width=0.5\textwidth]{figures/ch04/fig_IV_12.png}
\caption{}
\label{fig:IV-12}
\end{figure}

We can make a simple connection between the property of the
gradient just discussed and the fact that it is in the direction of the
greatest rate of increase.  Any displacement from the surface
$f(x,\,y,\,z) = \text{constant}$, regarded as a vector $\mathbf{s}$,
can be resolved into a component along the surface ($s_\parallel$) and
one normal to it ($s_\perp$), as shown in Figure~IV--12.  That part
of the displacement along the surface is ``wasted motion'' if our
aim in moving is to cause a change in the value of $f(x,\,y,\,z)$.
Only the normal component carries us away from the surface and
causes a change in $f$.  From this it is clear that the greatest
increase possible for a given magnitude of displacement should
occur when we move away from the surface in the normal direction.
But we have already established that the greatest rate of increase
occurs in the direction of the gradient.  Thus the gradient is normal
to the surface.


%% ================================================================
\section{The Gradient in Cylindrical and Spherical Coordinates}
%% ================================================================

A by-product of our discussion of the directional derivative is the
``easier and faster'' method for calculating the gradient in spherical
and cylindrical coordinates mentioned earlier (see page~120).  To
determine this method, we begin by outlining our derivation of
$df/ds$:\footnote{The calculation outlined here pertains to a function
of three variables and is a simple generalization of the calculation
on pages~130--133, which deals with a function of two variables.}

\begin{enumerate}
\item Our first step is to consider a scalar function of three
  Cartesian coordinates $f(x,\,y,\,z)$ and use Taylor series to
  determine the change in $f$ caused by a displacement from the
  point $(x,\,y,\,z)$ to a second point
  $(x+\Delta x,\,y+\Delta y,\,z+\Delta z)$.  We find for this change
  \[
  \Delta f = \pd{f}{x}\,\Delta x + \pd{f}{y}\,\Delta y
           + \pd{f}{z}\,\Delta z + \cdots\,.
  \]

\item We next write $\Delta f$ in terms of $\Delta\mathbf{s}$, the
  vector displacement from $(x,\,y,\,z)$ to
  $(x+\Delta x,\,y+\Delta y,\,z+\Delta z)$.  Clearly (see
  Figure~IV--13),
  \[
  \Delta\mathbf{s} = \bfi\,\Delta x + \bfj\,\Delta y + \bfk\,\Delta z,
  \]
  so that
  \[
  \Delta f = \left(\bfi\,\pd{f}{x} + \bfj\,\pd{f}{y}
            + \bfk\,\pd{f}{z}\right)\cdot\Delta\mathbf{s} + \cdots\,.
  \]

\begin{figure}[ht]
\centering
\includegraphics[width=0.5\textwidth]{figures/ch04/fig_IV_13.png}
\caption{}
\label{fig:IV-13}
\end{figure}

\item Finally, we write $\Delta\mathbf{s} = \hat{\mathbf{u}}\,\Delta s$,
  divide by $\Delta s$, and take the limit:
  \[
  \lim_{\Delta s\to 0}\frac{\Delta f}{\Delta s} \equiv \frac{df}{ds}
  = \left(\bfi\,\pd{f}{x} + \bfj\,\pd{f}{y}
    + \bfk\,\pd{f}{z}\right)\cdot\hat{\mathbf{u}}.
  \]
  The quantity that is dotted into $\hat{\mathbf{u}}$ in this last
  expression is then recognized as the gradient of $f$ in Cartesian
  coordinates.
\end{enumerate}

To obtain the gradient of a scalar function in cylindrical
coordinates we proceed in much the same way:

\begin{enumerate}
\item We consider a scalar function of three cylindrical coordinates,
  $f(r,\,\theta,\,z)$.  Using Taylor series, we find the change in
  $f$ due to a displacement from the point $(r,\,\theta,\,z)$ to a
  second point $(r+\Delta r,\,\theta+\Delta\theta,\,z+\Delta z)$:
  \[
  \Delta f = \pd{f}{r}\,\Delta r + \pd{f}{\theta}\,\Delta\theta
           + \pd{f}{z}\,\Delta z + \cdots\,.
  \]

\item Next, we write $\Delta f$ in terms of $\Delta\mathbf{s}$.  This
  is the heart of the calculation.  From Figure~IV--14 we have
  \[
  \Delta\mathbf{s}
  = \hat{\mathbf{e}}_r\,\Delta r
  + \hat{\mathbf{e}}_\theta\,r\,\Delta\theta
  + \hat{\mathbf{e}}_z\,\Delta z.
  \]

\begin{figure}[ht]
\centering
\includegraphics[width=0.5\textwidth]{figures/ch04/fig_IV_14.png}
\caption{}
\label{fig:IV-14}
\end{figure}

  Two features of this expression require some discussion.  First,
  the displacement in the direction of increasing $\theta$ (of
  magnitude $r\,\Delta\theta$) is an arc of a circle rather than a
  straight line segment.  However, since we will eventually pass to
  the limit as $\Delta s \to 0$, we may regard $\Delta\theta$ (as
  well as $\Delta r$ and $\Delta z$) as arbitrarily small, in which
  case the arc is arbitrarily close to its subtending chord.  Thus,
  as indicated in Figure~IV--15, $\Delta r$, $r\,\Delta\theta$, and
  $\Delta z$

\begin{figure}[ht]
\centering
\includegraphics[width=0.45\textwidth]{figures/ch04/fig_IV_15.png}
\caption{}
\label{fig:IV-15}
\end{figure}

  \noindent approximate to any desired degree of accuracy three
  mutually perpendicular displacements, the analogs of the three
  Cartesian displacements $\Delta x$, $\Delta y$, and $\Delta z$ (see
  Figure~IV--13).

  The second feature of our expression for $\Delta\mathbf{s}$ that
  requires comment also has to do with the displacement in the
  direction of increasing $\theta$.  It is this: Since the arc is part
  of a circle of radius $r + \Delta r$, we should, strictly speaking,
  write the displacement as $(r + \Delta r)\,\Delta\theta$, not
  $r\,\Delta\theta$.  But the additional term $\Delta r\,\Delta\theta$
  is ``second order''; that is, it is the product of two small
  quantities and therefore negligible compared with
  $r\,\Delta\theta$.

  Note the factor $1/r$ in the second term to compensate for the
  factor $r$ in $\hat{\mathbf{e}}_\theta\,r\,\Delta\theta$ in
  $\Delta\mathbf{s}$.

  If we now write our expression for $\Delta f$ in terms of
  $\Delta\mathbf{s}$, we get
  \[
  \Delta f = \left(\hat{\mathbf{e}}_r\,\pd{f}{r}
    + \hat{\mathbf{e}}_\theta\,\frac{1}{r}\,\pd{f}{\theta}
    + \hat{\mathbf{e}}_z\,\pd{f}{z}\right)
    \cdot\Delta\mathbf{s} + \cdots\,.
  \]

\item Finally, putting $\Delta\mathbf{s} = \hat{\mathbf{u}}\,\Delta s$,
  we find
  \[
  \lim_{\Delta s\to 0}\frac{\Delta f}{\Delta s} \equiv \frac{df}{ds}
  = \left(\hat{\mathbf{e}}_r\,\pd{f}{r}
    + \hat{\mathbf{e}}_\theta\,\frac{1}{r}\,\pd{f}{\theta}
    + \hat{\mathbf{e}}_z\,\pd{f}{z}\right)
    \cdot\hat{\mathbf{u}}.
  \]
  The quantity in the preceding expression dotted into
  $\hat{\mathbf{u}}$ is the gradient of $f$ in cylindrical coordinates.
\end{enumerate}

An analogous procedure can be used to find the gradient in
spherical coordinates; this has been left as an exercise (see
Problem~IV--22).


%% ================================================================
\section{Problems}
%% ================================================================

\begin{problems}

\prob
\begin{enumerate}[label=(\alph*)]
\item Calculate $\bfF = \nabla f$ for each of the following scalar functions:
  \begin{enumerate}[label=(\roman*)]
  \item $f = xyz$.
  \item $f = x^2 + y^2 + z^2$.
  \item $f = xy + yz + xz$.
  \item $f = 3x^2 - 4z^2$.
  \item $f = e^{-x}\sin y$.
  \end{enumerate}
\item Verify that
  \[
  \oint_C \bfF\cdot\hat{\mathbf{t}}\ds = 0
  \]
  for one or more of the functions $\bfF$ determined in part~(a)
  choosing for the curve $C$:
  \begin{enumerate}[label=(\roman*)]
  \item the square in the $xy$-plane with vertices at $(0,\,0)$,
    $(1,\,0)$, $(1,\,1)$, and $(0,\,1)$.
  \item the triangle in the $yz$-plane with vertices at $(0,\,0)$,
    $(1,\,0)$, and $(0,\,1)$.
  \item the circle of unit radius centered at the origin and lying in
    the $xz$-plane.
  \end{enumerate}
\item Verify by direct calculation that $\nabla\times\bfF = 0$ for one
  or more of the functions $\bfF$ determined in part~(a).
\end{enumerate}

\prob
Verify the following identities in which $f$ and $g$ are arbitrary
differentiable scalar functions of position, and $\bfF$ and $\bfG$
are arbitrary differentiable vector functions of position.
\begin{enumerate}[label=(\alph*)]
\item $\nabla(fg) = f\,\nabla g + g\,\nabla f$.
\item $\nabla(\bfF\cdot\bfG) = (\bfG\cdot\nabla)\bfF
  + (\bfF\cdot\nabla)\bfG + \bfF\times(\nabla\times\bfG)
  + \bfG\times(\nabla\times\bfF)$.
\item $\nabla\cdot(f\bfF) = f\,\nabla\cdot\bfF
  + \bfF\cdot\nabla f$.
\item $\nabla\cdot(\bfF\times\bfG) = \bfG\cdot(\nabla\times\bfF)
  - \bfF\cdot(\nabla\times\bfG)$.
\item $\nabla\times(f\bfF) = f\,\nabla\times\bfF
  + (\nabla f)\times\bfF$.
\item $\nabla\times(\bfF\times\bfG) = (\bfG\cdot\nabla)\bfF
  - (\bfF\cdot\nabla)\bfG + \bfF(\nabla\cdot\bfG)
  - \bfG(\nabla\cdot\bfF)$.
\item $\nabla\times(\nabla\times\bfF)
  = \nabla(\nabla\cdot\bfF) - \laplacian\bfF$.
\end{enumerate}

\prob
Show that $\nabla\times\nabla f = 0$ where $f(x,\,y,\,z)$ is an
arbitrary differentiable scalar function.  Assume that mixed
second-order partial derivatives are independent of the order of
differentiation.  For example,
$\partial^2 f/\partial x\,\partial z = \partial^2 f/\partial z\,\partial x$.

\prob
\begin{enumerate}[label=(\alph*)]
\item Each of the following functions is smooth in a simply connected
  region.  Determine which of them may be written as the gradient of a
  scalar function, and for those that can, use Equation~(\ref{eq:IV-2})
  to find that scalar function.
  \begin{enumerate}[label=(\roman*)]
  \item $\bfF = \bfi\,y$.
  \item $\bfF = C\bfk$, \quad $C$ a constant.
  \item $\bfF = \bfi\,yz + \bfj\,xz + \bfk\,xy$.
  \item $\bfF = \bfi\,x + \bfj\,y + \bfk\,z$.
  \item $\bfF = \bfi\,e^{-z}\sin y + \bfj\,e^{-y}\sin z
    + \bfk\,e^{-x}\sin y$.
  \end{enumerate}
\item Neither of the following functions is smooth everywhere.
  Nonetheless each can be written as the gradient of a scalar function.
  Use Equation~(\ref{eq:IV-2}) to find that scalar function.
  \begin{enumerate}[label=(\roman*)]
  \item $\bfF = \bfr/r^2$, \quad $\bfr = \bfi\,x + \bfj\,y$.
  \item $\bfF = \bfr/r^{1/2}$, \quad
    $\bfr = \bfi\,x + \bfj\,y + \bfk\,z$.
  \end{enumerate}
\end{enumerate}

\prob
The function $\bfF(r,\,\theta,\,z)$ defined in Problem~III--17 is
smooth and has zero curl in a nonsimply connected region consisting
of all of three-dimensional space with the $z$-axis removed.  Show
that there is no scalar function $\psi$ such that
$\bfF = \nabla\psi$ by evaluating the line integral of
$\bfF\cdot\hat{\mathbf{t}}$ from the point $P_1(0,\,-1,\,0)$ to the
point $P_2(0,\,1,\,0)$ over two different paths: $C_R$, the
right-hand side of the circle of radius~1 lying in the $xy$-plane and
centered at the origin (see figure), and $C_L$, the left-hand side of
the same circle.  Orient the paths as shown.  Why does the fact that
the two paths give different results imply that there is no scalar
function $\psi$ such that $\bfF = \nabla\psi$?

\begin{figure}[ht]
\centering
\includegraphics[width=0.4\textwidth]{figures/ch04/fig_IV_prob5.png}
\caption{}
\label{fig:IV-prob5}
\end{figure}

\prob
\begin{enumerate}[label=(\alph*)]
\item An electric dipole of strength $p$ situated at the origin and
  oriented in the positive $z$-direction gives rise to an
  electrostatic field
  \[
  \bfE(r,\,\theta,\,\phi)
  = \frac{1}{4\pi\epszero}\,\frac{p}{r^3}
    (2\hat{\mathbf{e}}_r\cos\phi + \hat{\mathbf{e}}_\phi\sin\phi)
  \]

\begin{figure}[ht]
\centering
\includegraphics[width=0.45\textwidth]{figures/ch04/fig_IV_prob6.png}
\caption{}
\label{fig:IV-prob6}
\end{figure}

  Use Equation~(\ref{eq:IV-2}) to show that the dipole potential is
  given by
  \[
  \Phi(r,\,\theta,\,\phi)
  = \frac{1}{4\pi\epszero}\,\frac{p\cos\phi}{r^2}\,.
  \]

  Useful information: In spherical coordinates,
  \[
  \hat{\mathbf{t}}
  = \hat{\mathbf{e}}_r\,\frac{dr}{ds}
  + \hat{\mathbf{e}}_\phi\,r\,\frac{d\phi}{ds}
  + \hat{\mathbf{e}}_\theta\,r\sin\phi\,\frac{d\theta}{ds}\,.
  \]

\item Calculate the flux of the dipole field through a sphere of
  radius $R$ centered at the origin.

\item What is the flux of the dipole field over \emph{any} closed
  surface that does not pass through the origin?
\end{enumerate}

\prob
Here is a ``proof'' that there is no such thing as magnetism.  One
of Maxwell's equations tells us that
\[
\nabla\cdot\bfB = 0,
\]
where $\bfB$ is any magnetic field.  Then using the divergence
theorem, we find
\[
\iint_S \bfB\cdot\nhat\dS = \iiint_V \nabla\cdot\bfB\dV = 0.
\]

Because $\bfB$ has zero divergence, we know (see Problem~III--24)
there exists a vector function, call it $\bfA$, such that
\[
\bfB = \nabla\times\bfA.
\]

Combining these last two equations, we get
\[
\iint_S \nhat\cdot\nabla\times\bfA\dS = 0.
\]

Next we apply Stokes' theorem and the preceding result to find
\[
\oint_C \bfA\cdot\hat{\mathbf{t}}\ds
= \iint_S \nhat\cdot\nabla\times\bfA\dS = 0.
\]

Thus we have shown that the circulation of $\bfA$ is path
independent.  It follows that we can write
$\bfA = \nabla\psi$, where $\psi$ is some scalar function.  Since
the curl of the gradient of a function is zero, we arrive at the
remarkable fact that
\[
\bfB = \nabla\times\nabla\psi = 0;
\]
that is, all magnetic fields are zero!  Where did we go wrong?
[Taken from G.~Arfken, \emph{Amer.\ J.\ Phys.}, \textbf{27}, 526
(1959).]

\prob
Fick's law states that in certain diffusion processes the current
density $\bfJ$ is proportional to the negative of the gradient of the
density $\rho$; that is, $\bfJ = -k\,\nabla\rho$, where $k$ is a
positive constant.  If a substance of density $\rho(x,\,y,\,z,\,t)$
and velocity $\mathbf{v}(x,\,y,\,z,\,t)$ diffuses according to Fick's
law, show that the flow is \emph{irrotational} (that is,
$\nabla\times\mathbf{v} = 0$).

\prob
\begin{enumerate}[label=(\alph*)]
\item A substance diffuses according to Fick's law (see
  Problem~IV--8).  Assuming the diffusing matter is conserved, derive
  the diffusion equation
  \[
  \pd{\rho}{t} = k\,\laplacian\rho.
  \]

\item Bacteria of density $\rho$ diffuse in a medium according to
  Fick's law and reproduce at a rate $\lambda\rho$ per unit volume
  ($\lambda$ is a positive constant).  Show that
  \[
  \pd{\rho}{t} = k\,\laplacian\rho + \lambda\rho.
  \]
\end{enumerate}

\prob
\begin{enumerate}[label=(\alph*)]
\item A fluid is said to be \emph{incompressible} if its density
  $\rho$ is a constant (that is, is independent of $x$, $y$, $z$, and
  $t$).  Use the continuity equation to show that the velocity
  $\mathbf{v}$ of an incompressible fluid satisfies the equation
  $\nabla\cdot\mathbf{v} = 0$.

\item If $\nabla\times\mathbf{v} = 0$, the fluid flow is said to be
  \emph{irrotational}.  Show that for an incompressible fluid
  undergoing irrotational flow,
  \[
  \laplacian\phi = 0,
  \]
  where $\phi$, a scalar function called the \emph{velocity
  potential}, is so defined that $\mathbf{v} = \nabla\phi$.
\end{enumerate}

\prob
The heat $Q$ in a body of volume $V$ is given by
\[
Q = c\iiint_V T\rho\dV,
\]
where $c$ is a constant called the specific heat of the body, and
$T(x,\,y,\,z,\,t)$ and $\rho(x,\,y,\,z)$ are, respectively, the
temperature and density of the body.  (Note that we are assuming the
density to be independent of time.)  The rate at which heat flows
through $S$, the bounding surface of the body, is given by
\[
\frac{dQ}{dt} = k\iint_S \nhat\cdot\nabla T\dS,
\]
where $k$ (assumed constant) is the thermal conductivity of the
body, and the integral is taken over the surface $S$ bounding the
body.  Use these facts to derive the heat flow equation
\[
\laplacian T = \alpha\,\pd{T}{t}\,,
\]
where $\alpha = c\rho/k$.

\prob
In nonrelativistic quantum mechanics a particle of mass $m$ moving
in a potential $V(x,\,y,\,z)$ is described by the Schr\"odinger
equation
\[
-\frac{\hbar^2}{2m}\,\laplacian\psi + V\psi
= i\hbar\,\pd{\psi}{t}\,,
\]
where $\hbar$ is Planck's constant divided by $2\pi$ and
$\psi(x,\,y,\,z,\,t)$, which is complex, is called the wave function.
The quantity $\rho = \psi^*\psi$ is interpreted as the probability
density.
\begin{enumerate}[label=(\alph*)]
\item Use the Schr\"odinger equation to derive an equation of the form
  \[
  \pd{\rho}{t} + \nabla\cdot\bfJ = 0
  \]
  and obtain thereby an expression for $\bfJ$ in terms of $\psi$,
  $\psi^*$, $m$, and $\hbar$.

\item Give an interpretation of $\bfJ$ and of the equation derived
  in~(a).
\end{enumerate}

\prob
\begin{enumerate}[label=(\alph*)]
\item Find the charge density $\rho(x,\,y,\,z)$ that produces the
  electric field
  \[
  \bfE = g(\bfi\,x + \bfj\,y + \bfk\,z),
  \]
  where $g$ is a constant.

\item Find an electrostatic potential $\Phi$ such that
  $-\nabla\Phi$ is the field $\bfE$ given in~(a).

\item Verify that $\laplacian\Phi = -\rho/\epszero$.
\end{enumerate}

\prob
\begin{enumerate}[label=(\alph*)]
\item Starting with the divergence theorem, derive the equation
  \[
  \iint_S \nhat\cdot(u\,\nabla v)\dS
  = \iiint_V [u\,\laplacian v + (\nabla u)\cdot(\nabla v)]\dV,
  \]
  where $u$ and $v$ are scalar functions of position and $S$ is a
  closed surface enclosing the volume $V$.  This is sometimes called
  the first form of Green's theorem.

\item If $\laplacian u = 0$ use the first form of Green's theorem to
  show that
  \[
  \iint_S \nhat\cdot(u\,\nabla u)\dS = \iiint_V |\nabla u|^2\dV,
  \]
  where $|\nabla u|^2 = (\nabla u)\cdot(\nabla u)$.

\item Use the first form of Green's theorem to show that
  \[
  \iint_S \nhat\cdot(u\,\nabla v - v\,\nabla u)\dS
  = \iiint_V (u\,\laplacian v - v\,\laplacian u)\dV.
  \]
  This is the second form of Green's theorem.
\end{enumerate}

\prob
An equation of the form
\[
\laplacian f = \frac{1}{v^2}\,\pdd{f}{t}\,,
\]
where $f$ is a twice-differentiable function of position and time, is
called a wave equation.  It describes a wave propagating in space
with velocity $v$.  Use Maxwell's equations (Problem~III--20) to
show that in the absence of charges and currents (that is, $\rho$ and
$\bfJ$ both zero), all three Cartesian components of both $\bfE$ and
$\bfB$ satisfy a wave equation with $v = c$, where
$c = 1/\sqrt{\epszero\mu_0}$ is the velocity of light.  For example,
\[
\laplacian E_x = \frac{1}{c^2}\,\pdd{E_x}{t}\,.
\]
Thus, the existence of electromagnetic waves traveling in empty
space with the velocity of light is a consequence of Maxwell's
equations.

\prob
\begin{enumerate}[label=(\alph*)]
\item In the text we found the potential and field for the case of an
  infinite cylinder between parallel plates with the cylinder held at
  zero potential.  How must the solution be modified if the cylinder is
  held at a potential $V_0 \neq 0$?

\item Show that there is no net charge on the cylinder.
\end{enumerate}

\prob
\begin{enumerate}[label=(\alph*)]
\item A sphere of radius $R$ is situated between two very large
  parallel plates that are separated by a distance $s$.  A potential
  difference is maintained between the plates and the sphere is held
  at zero potential.  Find the potential and field everywhere outside
  the sphere and between the plates.  Assume that $R \ll s$.

\item Show that there is no net charge on the sphere.

\item Repeat part~(a) assuming the sphere is held at a potential
  $V_0 \neq 0$.
\end{enumerate}

\prob
Let $f(x,\,y)$ be a differentiable scalar function of $x$ and $y$,
and let $\hat{\mathbf{u}} = \bfi\cos\theta + \bfj\sin\theta$.
Transform to a rotated coordinate system $x'$, $y'$ such that $x'$
is parallel to $\hat{\mathbf{u}}$ (see the figure).  Show that the
directional derivative in the direction of $\hat{\mathbf{u}}$ is given
by
\[
\frac{df}{ds} = \hat{\mathbf{u}}\cdot\nabla f = \pd{f}{x'}\,.
\]

\begin{figure}[ht]
\centering
\includegraphics[width=0.45\textwidth]{figures/ch04/fig_IV_prob18.png}
\caption{}
\label{fig:IV-prob18}
\end{figure}

\prob
You are at a point $(a,\,b,\,c)$ on the surface
\[
z = (r^2 - x^2 - y^2)^{1/2}\qquad (z \geq 0).
\]
Assuming both $a$ and $b$ are positive, in what direction must you
move
\begin{enumerate}[label=(\alph*)]
\item so that the rate of change of $z$ will be zero?
\item so that the rate of increase of $z$ will be greatest?
\item so that the rate of decrease of $z$ will be greatest?
\end{enumerate}
Draw a sketch to show the geometric significance of your answers.

\prob
The unit vector normal to the surface $z = f(x,\,y)$ is given by
\[
\nhat = \left(-\bfi\,\pd{f}{x} - \bfj\,\pd{f}{y} + \bfk\right)
\bigg/ \sqrt{1 + \left(\pd{f}{x}\right)^{\!2}
             + \left(\pd{f}{y}\right)^{\!2}}
\]
[see Equation~(II--4)].  We have also established that $\nabla F$ is
a vector normal to the surface $F(x,\,y,\,z) = \text{const.}$
(page~140) so that $\nabla F/|\nabla F|$ is a unit vector normal to
the surface $F(x,\,y,\,z) = \text{const}$.  Show that these two
expressions for the unit normal vector are identical if
$F(x,y,z) = \text{const.}$ and $z = f(x,\,y)$ describe the same
surface.

\prob
Use the results of Problem~II--18 and the expression for the
gradient in cylindrical coordinates (see page~144) to obtain the form
of the Laplacian in cylindrical coordinates given on page~129.

\prob
Using the procedure outlined in the text (pages~141--144) obtain the
expression for the gradient of $\psi$ in spherical coordinates:
\[
\nabla\psi
= \hat{\mathbf{e}}_r\,\pd{\psi}{r}
+ \hat{\mathbf{e}}_\phi\,\frac{1}{r}\,\pd{\psi}{\phi}
+ \hat{\mathbf{e}}_\theta\,\frac{1}{r\sin\phi}\,\pd{\psi}{\theta}\,.
\]

\prob
Use the results of Problem~II--19 and the expression for the
gradient in spherical coordinates derived in Problem~IV--22 to obtain
the form of the Laplacian in spherical coordinates given on
page~126.

\prob
Suppose you find a solution of Laplace's equation that satisfies
certain boundary conditions.  Is this solution unique or are there
others?  This problem will answer that question in certain simple
cases.  Consider the region of space completely enclosed by a
surface $S_0$ and containing in its interior objects $1, 2, 3, \ldots$
(two of which are pictured in the diagram).  Suppose that $S_0$ is
maintained at a constant potential $\Phi_0$, object no.~1 at
$\Phi_1$, object no.~2 at $\Phi_2$, and so on.  Then in the
charge-free region $R$ enclosed by $S_0$ and between the objects,
the potential must satisfy Laplace's equation
\[
\laplacian\Phi = 0
\]
and the boundary conditions
\[
\Phi = \begin{cases}
\Phi_0 & \text{on } S_0 \\
\Phi_1 & \text{on } S_1 \\
\Phi_2 & \text{on } S_2 \\
\vdots
\end{cases}
\]

\begin{figure}[ht]
\centering
\includegraphics[width=0.5\textwidth]{figures/ch04/fig_IV_prob24.png}
\caption{}
\label{fig:IV-prob24}
\end{figure}

The following steps will guide you through a proof that $\Phi$ is
unique.
\begin{enumerate}[label=(\alph*)]
\item Assume that there are two potentials $u$ and $v$, both of which
  satisfy Laplace's equation and the boundary conditions listed
  earlier.  Form their difference $w = u - v$.  Show that
  $\laplacian w = 0$ in $R$.

\item What are the boundary conditions satisfied by $w$?

\item Apply the divergence theorem to
  \[
  \iint_S \nhat\cdot(w\,\nabla w)\dS
  \]
  where the integration is carried out over the surface
  $S_0 + S_1 + S_2 + \cdots$, and show thereby that
  \[
  \iiint_V |\nabla w|^2\dV = 0,
  \]
  where $V$ is the volume of the region $R$.

\item From the result of~(c) argue that $\nabla w = 0$ and that this,
  in turn, means $w$ is a constant.

\item If $w$ is a constant, what is its value?  (Use the boundary
  conditions on $w$ to answer this.)  What does this say about $u$
  and $v$?

\item The uniqueness proof outlined in~(a) to~(e) involves specifying
  the value of the potential on various surfaces.  Might we have
  specified a different kind of boundary condition and still proved
  uniqueness?  If so, in what way or ways would the proof and the
  result differ from those given above?
\end{enumerate}

\prob
In the text we defined the gradient in terms of certain partial
derivatives.  It is possible to give an alternative definition similar
in form to our definitions of the divergence and the curl.  Thus,
\[
\nabla f = \lim_{\Delta V\to 0}\frac{1}{\Delta V}
  \iint_S \nhat\,f\dS.
\]
Here $f$ is a scalar function of position, $S$ a closed surface, and
$\Delta V$ the volume it encloses.  As usual, $\nhat$ is a unit vector
normal to $S$ and pointing out from the enclosed volume.
\begin{enumerate}[label=(\alph*)]
\item Following a procedure similar to the one used in the text in
  treating the divergence, integrate over a ``cuboid'' and show that
  the preceding definition yields the expression
  \[
  \nabla f = \bfi\,\pd{f}{x} + \bfj\,\pd{f}{y} + \bfk\,\pd{f}{z}\,.
  \]

\item Use the alternative definition of the gradient given above to
  show that the directional derivative of $f$ in the direction
  specified by the unit vector $\hat{\mathbf{u}}$ is given by
  \[
  \frac{df}{ds} = \hat{\mathbf{u}}\cdot\nabla f.
  \]
  [\emph{Hint:} Evaluate
  \[
  \hat{\mathbf{u}}\cdot\iint_S \nhat\,f\dS
  = \iint_S \hat{\mathbf{u}}\cdot\nhat\,f\dS
  \]
  over a small cylinder (length $\Delta s$, cross-sectional area
  $\Delta A$; see figure) whose axis is in the direction of the
  constant unit vector $\hat{\mathbf{u}}$.  Then divide by the volume
  of the cylinder ($\Delta s\,\Delta A$) and take the limit as the
  volume approaches zero.]

\begin{figure}[ht]
\centering
\includegraphics[width=0.45\textwidth]{figures/ch04/fig_IV_prob25b.png}
\caption{}
\label{fig:IV-prob25b}
\end{figure}

\item Arguing as we did in the text in establishing the divergence
  theorem, use the alternative definition of the gradient to show that
  \[
  \iint_S \nhat\,f\dS = \iiint_V \nabla f\dV,
  \]
  where $S$ is a closed surface enclosing the volume $V$.

\item Obtain the relation stated in~(c) directly from the divergence
  theorem.  [\emph{Hint:} In
  $\iint_S \bfF\cdot\nhat\dS = \iiint_V \nabla\cdot\bfF\dV$ put
  $\bfF = \hat{\mathbf{e}}\,f$ where $\hat{\mathbf{e}}$ is a constant
  unit vector.]

\item Verify the relation stated in~(c) for the scalar function
  \[
  f = x^2 + y^2 + z^2
  \]
  integrating over the unit cylinder shown in the figure.

\begin{figure}[ht]
\centering
\includegraphics[width=0.45\textwidth]{figures/ch04/fig_IV_prob25e.png}
\caption{}
\label{fig:IV-prob25e}
\end{figure}
\end{enumerate}

\prob
\begin{enumerate}[label=(\alph*)]
\item Consider a surface $z = f(x,\,y)$.  Let $\mathbf{u}$ be a vector
  of arbitrary length tangent to the surface at a point $P(x,\,y,\,z)$
  in the direction of the unit vector
  $\hat{\mathbf{p}} = \bfi\,p_x + \bfj\,p_y$ as indicated in the
  figure.  Use the directional derivative to show that
  \[
  \mathbf{u} = \hat{\mathbf{p}}
    + \bfk(\hat{\mathbf{p}}\cdot\nabla f),
  \]
  where $\nabla f$ is evaluated at $(x,\,y)$.  [\emph{Note:} Since
  the length of $\mathbf{u}$ is arbitrary, your result may differ from
  the preceding by some positive multiplicative constant.]

\item Let $\mathbf{v}$ be a second vector of arbitrary length tangent
  to the surface at $P$ but in the direction of the unit vector
  $\hat{\mathbf{q}} = \bfi\,q_x + \bfj\,q_y$
  ($\hat{\mathbf{p}} \neq \hat{\mathbf{q}}$).  Then from~(a) we have
  \[
  \mathbf{v} = \hat{\mathbf{q}}
    + \bfk(\hat{\mathbf{q}}\cdot\nabla f).
  \]
  Show that
  \[
  \mathbf{u}\times\mathbf{v}
  = [\bfk\cdot(\hat{\mathbf{p}}\times\hat{\mathbf{q}})]
    (\bfk - \nabla f)
  \]
  and use this to rederive Equation~(II--4) for the unit vector
  $\nhat$ normal to the surface $z = f(x,\,y)$ at $(x,\,y,\,z)$.
  This shows that the result derived in the text for $\nhat$ is unique
  (apart from sign) even though it was obtained with the special
  choices $\hat{\mathbf{p}} = \bfi$ and $\hat{\mathbf{q}} = \bfj$.

\begin{figure}[ht]
\centering
\includegraphics[width=0.5\textwidth]{figures/ch04/fig_IV_prob26.png}
\caption{}
\label{fig:IV-prob26}
\end{figure}
\end{enumerate}

\prob
\begin{enumerate}[label=(\alph*)]
\item Using Maxwell's equations (see Problem~III--20), show that we
  can write
  \[
  \bfB = \nabla\times\bfA,
  \]
  \[
  \bfE = -\nabla\Phi - \pd{\bfA}{t}\,,
  \]
  where $\bfA$ (called the vector potential) is some vector function of
  position and time, and $\Phi$ (the scalar potential) is some scalar
  function of position and time, provided $\bfA$ and $\Phi$ satisfy
  the equations
  \[
  \laplacian\Phi + \pd{}{t}(\nabla\cdot\bfA)
  = -\rho/\epszero,
  \]
  \[
  \laplacian\bfA - \mu_0\epszero\,\pdd{\bfA}{t}
  + \nabla\!\left[\nabla\cdot\bfA
    + \mu_0\epszero\,\pd{\Phi}{t}\right]
  = -\mu_0\bfJ.
  \]

\item Show that if we define two new potentials
  \[
  \bfA' = \bfA + \nabla\chi,\qquad
  \Phi' = \Phi - \pd{\chi}{t}\,,
  \]
  where $\chi$ is an arbitrary scalar function of position and time,
  then
  \[
  \bfB = \nabla\times\bfA',\qquad
  \bfE = -\nabla\Phi' - \pd{\bfA'}{t}\,.
  \]
  That is, the fields $\bfE$ and $\bfB$ are not modified by the change
  in the potentials $\bfA$ and $\Phi$.  The change from
  $(\bfA,\,\Phi)$ to $(\bfA',\,\Phi')$ is called a \emph{gauge
  transformation}.

\item Show that if we require $\chi$ to satisfy the equation
  \[
  \laplacian\chi - \epszero\mu_0\,\pdd{\chi}{t}
  = -\left[\nabla\cdot\bfA
    + \epszero\mu_0\,\pd{\Phi}{t}\right],
  \]
  then
  \[
  \nabla\cdot\bfA' + \epszero\mu_0\,\pd{\Phi'}{t} = 0.
  \]

\item If $\chi$ satisfies the equation given in~(c), show that $\bfA'$
  and $\Phi'$ satisfy the equations
  \[
  \laplacian\Phi' - \epszero\mu_0\,\pdd{\Phi'}{t}
  = -\frac{\rho}{\epszero}
  \]
  and
  \[
  \laplacian\bfA' - \epszero\mu_0\,\pdd{\bfA'}{t}
  = -\mu_0\bfJ.
  \]
  The point to all of this is that we can make a gauge transformation
  [as in~(b)], impose the condition given in~(c), and thereby obtain a
  scalar and a vector potential that satisfy the equations in part~(d),
  which are wave equations with source terms proportional to $\rho$
  and $\bfJ$.
\end{enumerate}

\prob
The equation of motion of an ideal fluid can be written
\[
\iiint_V \rho\,\mathbf{f}_{\text{ext}}\dV
- \iint_S \nhat\,p\dS
= \iiint_V \rho\left[\pd{\mathbf{v}}{t}
  + (\mathbf{v}\cdot\nabla)\mathbf{v}\right]\dV
\]
where $V$ is the volume of the fluid and $S$ is its surface.  Here
$\mathbf{f}_{\text{ext}}(x,\,y,\,z)$ is the external force per unit mass
acting on the fluid, $p(x,\,y,\,z)$ is the pressure of the fluid, and
$\rho(x,\,y,\,z)$ is its density, all at a point $(x,\,y,\,z)$ in the
fluid, and $\mathbf{v}(x,\,y,\,z,\,t)$ is the velocity of the fluid at
the point $(x,\,y,\,z)$ and at time $t$.
\begin{enumerate}[label=(\alph*)]
\item Use the form of the divergence theorem given in
  Problem~IV--25(c) to rewrite the equation of motion of an ideal
  fluid in the form
  \[
  \mathbf{f}_{\text{ext}} - \frac{1}{\rho}\,\nabla p
  = \pd{\mathbf{v}}{t} + (\mathbf{v}\cdot\nabla)\mathbf{v}.
  \]

\item Show that in the static case ($\mathbf{v} = 0$), the equation
  of motion becomes
  \[
  \mathbf{f}_{\text{ext}} = (1/\rho)\,\nabla p.
  \]

\item Consider a column of incompressible fluid oriented vertically
  parallel to the $z$-axis as shown in the figure.  Assuming that the
  only external force acting on the fluid is the downward uniform
  gravitational attraction of the earth, apply the equation for the
  static case given in~(b) to show that
  \[
  p = p_0 - \rho g z
  \]
  where $g$ is the acceleration due to gravity and $p_0$ is a
  constant.

\begin{figure}[ht]
\centering
\includegraphics[width=0.3\textwidth]{figures/ch04/fig_IV_prob28.png}
\caption{}
\label{fig:IV-prob28}
\end{figure}
\end{enumerate}

\end{problems}
